\chapter{Bạn Cùng Bàn Và Đổi Chỗ}
\chapterintro{Ngồi cạnh ai cũng là một loại số phận học đường}{Bạn cùng bàn có thể cứu việc học, cũng có thể phá nó rất duyên.}

\begin{lucbatbox}{Lục bát: Bạn cùng bàn nói nhiều}
Ngồi cạnh bạn quá hoạt ngôn
Ngày nào lớp cũng như còn giờ chơi
Bạn em kể chuyện không ngơi
Từ nhà đến lớp từ trời tới mưa
Bài thì thầy vẫn đang đưa
Riêng tai em đã say sưa chuyện người
\end{lucbatbox}

\begin{tutuyetbox}{Thất ngôn tứ tuyệt: Đổi chỗ vì học hành}
Em xin đổi chỗ rất ngay tình
Bảo rằng ngồi đó khó học ghê
Thật ra lý do không sâu lắm
Bạn kia nói chuyện suốt giờ nghe
\end{tutuyetbox}

\begin{ghichu}
Bạn cùng bàn là môn học không có trong thời khóa biểu nhưng ảnh hưởng sâu rộng tới kết quả học tập.
\end{ghichu}

\ngancach

\begin{lucbatbox}{Lục bát: Ngồi gần bạn quá giỏi}
Ngồi gần bạn học rất siêng
Em bỗng thấy mình nghiêng hẳn về lo
Bạn ghi chép kín từng tờ
Em còn chưa hiểu thầy vừa nói chi
Ngó sang càng thấy sầu bi
Cùng một chiếc ghế, hai khi khác đường
\end{lucbatbox}

\begin{tutuyetbox}{Thất ngôn tứ tuyệt: Bạn vui quá tay}
Bạn em tốt bụng hiền khô lắm
Có bánh chia đều có đồ chung
Chỉ khi vào đúng giờ học ấy
Bạn vui hơi quá cả bàn rung
\end{tutuyetbox}